\documentclass[11pt]{article}
\usepackage{amsmath, amssymb, geometry, graphicx, xcolor}
\geometry{a4paper, margin=1in}

\title{\textbf{Homework 15.7: Triple Integrals in Cylindrical Coordinates}}
\author{Tashfeen Omran}
\date{April 2024}

\begin{document}

\maketitle
\tableofcontents
\newpage

\section{Part 1: Comprehensive Introduction, Context, and Prerequisites}

\subsection{Core Concepts}
Triple integrals allow us to sum infinitely many infinitesimally small quantities over a three-dimensional region. While computing these in rectangular coordinates $(x, y, z)$ is often straightforward for box-like regions, doing so for regions exhibiting rotational symmetry (like cylinders, cones, and spheres) leads to complicated algebraic boundaries involving square roots. 

Cylindrical coordinates solve this by representing 3D space as a natural extension of 2D polar coordinates. Instead of $(x, y, z)$, we use $(r, \theta, z)$:
\begin{itemize}
    \item $r$ is the directed distance from the origin to the projection of the point onto the $xy$-plane.
    \item $\theta$ is the angle measured counterclockwise from the positive $x$-axis to the projection line.
    \item $z$ is the exact same vertical height as in rectangular coordinates.
\end{itemize}
When calculating a volume or a triple integral in this system, the infinitesimal volume element $dV = dx \, dy \, dz$ must be transformed. Because a small change in angle $\Delta\theta$ creates an arc length of $r \Delta\theta$, the area of the polar base is $r \, dr \, d\theta$. Extending this vertically by $dz$ yields the volume differential:
\[ dV = r \, dz \, dr \, d\theta \]
This extra $r$ (the Jacobian determinant of the transformation) is critical and represents how area scales with distance from the origin.

\subsection{Intuition and Derivation}
Imagine a standard soup can. Describing its boundary in $(x,y,z)$ requires the equation $x^2 + y^2 = R^2$, which forces us to use integrals with bounds like $-\sqrt{R^2 - x^2}$ to $\sqrt{R^2 - x^2}$. By switching to cylindrical coordinates, the boundary of the can is simply $r = R$. A complicated curve becomes a constant! The integration bounds become purely numerical, turning a tedious calculus problem into a simple arithmetic one. The derivation relies on basic trigonometry for the $xy$-plane (polar coordinates) while leaving the $z$-axis untouched.

\subsection{Historical Context and Motivation}
The formalization of multiple coordinate systems arose in the 18th and 19th centuries through the work of Euler, Lagrange, and Gauss. Mathematicians and physicists were trying to model fluid dynamics, heat distribution, and gravitational pull for celestial bodies. Rectangular coordinates made the differential equations for these symmetrical, rounded shapes practically impossible to solve. Cylindrical and spherical coordinates were developed to align the mathematical grid with the natural symmetry of the physical objects, drastically simplifying the resulting equations. 

\subsection{Key Formulas}
\textbf{Coordinate Conversions:}
\[ x = r \cos(\theta) \quad \text{and} \quad y = r \sin(\theta) \]
\[ r^2 = x^2 + y^2 \quad \text{and} \quad \tan(\theta) = \frac{y}{x} \]
\[ z = z \]

\textbf{Triple Integral in Cylindrical Coordinates:}
\[ \iiint_E f(x, y, z) \, dV = \int_{\alpha}^{\beta} \int_{h_1(\theta)}^{h_2(\theta)} \int_{u_1(r, \theta)}^{u_2(r, \theta)} f(r\cos(\theta), r\sin(\theta), z) \, r \, dz \, dr \, d\theta \]

\subsection{Prerequisites}
To succeed in this topic, you must master:
\begin{enumerate}
    \item \textbf{Polar Coordinates:} Fluency in transforming $(x,y)$ to $(r,\theta)$ and understanding polar curves.
    \item \textbf{Iterated Integrals:} Evaluating double and triple integrals step-by-step from the inside out.
    \item \textbf{Trigonometric Integration:} Using identities like $\cos^2(\theta) = \frac{1 + \cos(2\theta)}{2}$ and computing basic integrals of sine and cosine.
    \item \textbf{Algebraic Substitution (u-sub):} Vital for integrating terms that result from the extra Jacobian $r$.
\end{enumerate}

\newpage
\section{Part 2: Detailed Homework Solutions}

\subsection*{Problem 1}
\textbf{Question:} Plot the point whose cylindrical coordinates are given. Then find the rectangular coordinates of the point. (a) $(6, \frac{\pi}{2}, 3)$, (b) $(8, -\frac{\pi}{4}, -7)$.

\textbf{Solution (a):}
We are given $(r, \theta, z) = \left(6, \frac{\pi}{2}, 3\right)$.
To find $(x, y, z)$:
\[ x = r \cos(\theta) = 6 \cos\left(\frac{\pi}{2}\right) = 6(0) = 0 \]
\[ y = r \sin(\theta) = 6 \sin\left(\frac{\pi}{2}\right) = 6(1) = 6 \]
The $z$-coordinate remains the same: $z = 3$.
\textbf{Final Answer (a):} $(0, 6, 3)$

\textbf{Solution (b):}
We are given $(r, \theta, z) = \left(8, -\frac{\pi}{4}, -7\right)$.
\[ x = r \cos(\theta) = 8 \cos\left(-\frac{\pi}{4}\right) = 8 \left(\frac{\sqrt{2}}{2}\right) = 4\sqrt{2} \]
\[ y = r \sin(\theta) = 8 \sin\left(-\frac{\pi}{4}\right) = 8 \left(-\frac{\sqrt{2}}{2}\right) = -4\sqrt{2} \]
The $z$-coordinate remains: $z = -7$.
\textbf{Final Answer (b):} $(4\sqrt{2}, -4\sqrt{2}, -7)$

---

\subsection*{Problem 2}
\textbf{Question:} Plot the point whose cylindrical coordinates are given. Then find the rectangular coordinates. (a) $(6, \frac{5\pi}{6}, 2)$, (b) $(6, -\frac{2\pi}{3}, 5)$.

\textbf{Solution (a):}
Given $(r, \theta, z) = \left(6, \frac{5\pi}{6}, 2\right)$.
\[ x = 6 \cos\left(\frac{5\pi}{6}\right) = 6 \left(-\frac{\sqrt{3}}{2}\right) = -3\sqrt{3} \]
\[ y = 6 \sin\left(\frac{5\pi}{6}\right) = 6 \left(\frac{1}{2}\right) = 3 \]
\textbf{Final Answer (a):} $(-3\sqrt{3}, 3, 2)$

\textbf{Solution (b):}
Given $(r, \theta, z) = \left(6, -\frac{2\pi}{3}, 5\right)$.
\[ x = 6 \cos\left(-\frac{2\pi}{3}\right) = 6 \left(-\frac{1}{2}\right) = -3 \]
\[ y = 6 \sin\left(-\frac{2\pi}{3}\right) = 6 \left(-\frac{\sqrt{3}}{2}\right) = -3\sqrt{3} \]
\textbf{Final Answer (b):} $(-3, -3\sqrt{3}, 5)$

---

\subsection*{Problem 3}
\textbf{Question:} Change from rectangular to cylindrical coordinates. (a) $(3, 3, -8)$, (b) $(4\sqrt{3}, -4, \sqrt{3})$.

\textbf{Solution (a):}
Given $(x, y, z) = (3, 3, -8)$.
\[ r = \sqrt{x^2 + y^2} = \sqrt{3^2 + 3^2} = \sqrt{18} = 3\sqrt{2} \]
\[ \tan(\theta) = \frac{y}{x} = \frac{3}{3} = 1 \]
Since $x > 0$ and $y > 0$, the point is in the first quadrant. Thus, $\theta = \frac{\pi}{4}$.
The $z$-coordinate is unchanged.
\textbf{Final Answer (a):} $\left(3\sqrt{2}, \frac{\pi}{4}, -8\right)$

\textbf{Solution (b):}
Given $(x, y, z) = (4\sqrt{3}, -4, \sqrt{3})$.
\[ r = \sqrt{(4\sqrt{3})^2 + (-4)^2} = \sqrt{48 + 16} = \sqrt{64} = 8 \]
\[ \tan(\theta) = \frac{y}{x} = \frac{-4}{4\sqrt{3}} = -\frac{1}{\sqrt{3}} \]
Since $x > 0$ and $y < 0$, the point is in the fourth quadrant. The standard angle is $\theta = \frac{11\pi}{6}$ (or $-\frac{\pi}{6}$).
\textbf{Final Answer (b):} $\left(8, \frac{11\pi}{6}, \sqrt{3}\right)$

---

\subsection*{Problem 4}
\textbf{Question:} Identify the surface whose equation is given: $r^2 + z^2 = 64$.
This is a [Select] with center the [Select] and radius [Box].

\textbf{Solution:}
We know the conversion formula $r^2 = x^2 + y^2$. Substituting this into the equation:
\[ x^2 + y^2 + z^2 = 64 \]
This is the standard equation of a sphere centered at the origin $(0,0,0)$ with a radius of $\sqrt{64} = 8$.
\textbf{Final Answer:} This is a \textbf{sphere} with center the \textbf{origin} and radius \textbf{8}.

---

\subsection*{Problem 5}
\textbf{Question:} Write the equations in cylindrical coordinates. (a) $5x^2 - 6x + 5y^2 + z^2 = 5$, (b) $z = 5x^2 - 5y^2$.

\textbf{Solution (a):}
Group the terms: $5(x^2 + y^2) - 6x + z^2 = 5$.
Substitute $x^2 + y^2 = r^2$ and $x = r\cos(\theta)$:
\[ 5r^2 - 6(r\cos(\theta)) + z^2 = 5 \]
\textbf{Final Answer (a):} $5r^2 - 6r\cos(\theta) + z^2 = 5$

\textbf{Solution (b):}
Substitute $x = r\cos(\theta)$ and $y = r\sin(\theta)$:
\[ z = 5(r\cos(\theta))^2 - 5(r\sin(\theta))^2 = 5r^2\cos^2(\theta) - 5r^2\sin^2(\theta) \]
Factor out $5r^2$:
\[ z = 5r^2(\cos^2(\theta) - \sin^2(\theta)) \]
Using the double angle identity $\cos(2\theta) = \cos^2(\theta) - \sin^2(\theta)$:
\textbf{Final Answer (b):} $z = 5r^2\cos(2\theta)$

---

\subsection*{Problem 6}
\textbf{Question:} Write the equations in cylindrical coordinates. (a) $4x^2 + 4y^2 - z^2 = 9$, (b) $7x - y + z = 3$.

\textbf{Solution (a):}
Factor out the 4: $4(x^2 + y^2) - z^2 = 9$.
Substitute $r^2 = x^2 + y^2$:
\textbf{Final Answer (a):} $4r^2 - z^2 = 9$

\textbf{Solution (b):}
Substitute $x = r\cos(\theta)$ and $y = r\sin(\theta)$:
\[ 7(r\cos(\theta)) - (r\sin(\theta)) + z = 3 \]
\textbf{Final Answer (b):} $7r\cos(\theta) - r\sin(\theta) + z = 3$

---

\subsection*{Problem 7}
\textbf{Question:} $f(x, y, z) = 8(x^2 + y^2)$. The solid $E$ is bounded by $z = 2 - x^2 - y^2$ and $z = x^2 + y^2$.
(a) Express the triple integral $\iiint_E f(x, y, z) \, dV$ as an iterated integral in cylindrical coordinates.
(b) Evaluate the iterated integral.

\textbf{Solution (a):}
First, convert the boundary equations to cylindrical coordinates:
Upper boundary: $z = 2 - r^2$
Lower boundary: $z = r^2$
Find the intersection to get the radial bounds by setting them equal:
\[ 2 - r^2 = r^2 \implies 2 = 2r^2 \implies r^2 = 1 \implies r = 1 \]
So, $0 \leq r \leq 1$. The region completes a full circle, so $0 \leq \theta \leq 2\pi$.
The integrand is $f(r, \theta, z) = 8(r^2)$.
The volume differential is $dV = r \, dz \, dr \, d\theta$.
The integral setup is:
\[ \int_{0}^{2\pi} \int_{0}^{1} \int_{r^2}^{2-r^2} (8r^2) \cdot r \, dz \, dr \, d\theta = \int_{0}^{2\pi} \int_{0}^{1} \int_{r^2}^{2-r^2} 8r^3 \, dz \, dr \, d\theta \]
\textbf{Final Answer (a):} $\int_{0}^{2\pi} \int_{0}^{1} \int_{r^2}^{2-r^2} 8r^3 \, dz \, dr \, d\theta$

\textbf{Solution (b):}
Evaluate the inner integral with respect to $z$:
\[ \int_{r^2}^{2-r^2} 8r^3 \, dz = 8r^3 [z]_{r^2}^{2-r^2} = 8r^3 ((2 - r^2) - r^2) = 8r^3(2 - 2r^2) = 16r^3 - 16r^5 \]
Evaluate the middle integral with respect to $r$:
\[ \int_{0}^{1} (16r^3 - 16r^5) \, dr = \left[ 4r^4 - \frac{16}{6}r^6 \right]_{0}^{1} = \left(4 - \frac{8}{3}\right) - (0) = \frac{12}{3} - \frac{8}{3} = \frac{4}{3} \]
Evaluate the outer integral with respect to $\theta$:
\[ \int_{0}^{2\pi} \frac{4}{3} \, d\theta = \frac{4}{3} (2\pi) = \frac{8\pi}{3} \]
\textbf{Final Answer (b):} $\frac{8\pi}{3}$

---

\subsection*{Problem 8}
\textbf{Question:} Sketch the solid whose volume is given by the integral and evaluate it: $\int_0^5 \int_0^{2\pi} \int_0^r r \, dz \, d\theta \, dr$.

\textbf{Solution:}
First, interpret the bounds:
$z$ goes from $0$ (the $xy$-plane) to $z = r$. The equation $z = r$ represents a cone opening upwards.
$\theta$ goes from $0$ to $2\pi$, representing a full rotation.
$r$ goes from $0$ to $5$, meaning the solid is restricted within a cylinder of radius $5$.
Therefore, the solid lies above the $xy$-plane, below the cone $z = \sqrt{x^2 + y^2}$, and inside the cylinder $x^2 + y^2 = 25$.

Now, let's evaluate the integral:
Inner integral:
\[ \int_0^r r \, dz = r [z]_0^r = r(r - 0) = r^2 \]
Middle integral:
\[ \int_0^{2\pi} r^2 \, d\theta = r^2 [\theta]_0^{2\pi} = 2\pi r^2 \]
Outer integral:
\[ \int_0^5 2\pi r^2 \, dr = 2\pi \left[ \frac{r^3}{3} \right]_0^5 = 2\pi \left( \frac{125}{3} \right) = \frac{250\pi}{3} \]
\textbf{Final Answer:} The solid lies inside the cylinder $r=5$, above $z=0$, and below the cone $z=r$. The evaluated value is $\frac{250\pi}{3}$.

---

\subsection*{Problem 9}
\textbf{Question:} Evaluate $\iiint_E \sqrt{x^2 + y^2} \, dV$, where $E$ is the region that lies inside the cylinder $x^2 + y^2 = 4$ and between the planes $z = 2$ and $z = 10$.

\textbf{Solution:}
Convert to cylindrical coordinates:
The cylinder $x^2 + y^2 = 4$ becomes $r^2 = 4 \implies r = 2$. Limits for $r$ are $0 \leq r \leq 2$.
Limits for $\theta$ are $0 \leq \theta \leq 2\pi$.
Limits for $z$ are given as $2 \leq z \leq 10$.
The integrand $\sqrt{x^2 + y^2}$ becomes $r$.
The volume element $dV = r \, dz \, dr \, d\theta$.
The integral is:
\[ \int_0^{2\pi} \int_0^2 \int_2^{10} (r)(r) \, dz \, dr \, d\theta = \int_0^{2\pi} \int_0^2 \int_2^{10} r^2 \, dz \, dr \, d\theta \]
Evaluate the inner integral:
\[ \int_2^{10} r^2 \, dz = r^2 [z]_2^{10} = r^2 (10 - 2) = 8r^2 \]
Evaluate the middle integral:
\[ \int_0^2 8r^2 \, dr = \left[ \frac{8r^3}{3} \right]_0^2 = \frac{8(8)}{3} = \frac{64}{3} \]
Evaluate the outer integral:
\[ \int_0^{2\pi} \frac{64}{3} \, d\theta = \frac{64}{3} [\theta]_0^{2\pi} = \frac{128\pi}{3} \]
\textbf{Final Answer:} $\frac{128\pi}{3}$

---

\subsection*{Problem 10}
\textbf{Question:} Consider $\iiint_E z \, dV$, where $E$ is enclosed by the paraboloid $z = x^2 + y^2$ and the plane $z = 4$. Write the above using cylindrical coordinates. Choose $0 < A \leq 2\pi$. Find $A, B, C$, and evaluate.
Integral format: $\int_0^A \int_0^B \int_C^4 (\dots) \, dz \, dr \, d\theta$.

\textbf{Solution:}
Boundaries in cylindrical coordinates:
Bottom surface: $z = r^2$. (This is the lower bound $C$).
Top surface: $z = 4$.
Intersection to find $r$-bounds: $r^2 = 4 \implies r = 2$. So $0 \leq r \leq 2$. This gives $B = 2$.
The region is rotationally symmetric, so $0 \leq \theta \leq 2\pi$. This gives $A = 2\pi$.
The integrand is $z$. The differential is $r \, dz \, dr \, d\theta$.
The integral setup is:
\[ \int_0^{2\pi} \int_0^2 \int_{r^2}^4 zr \, dz \, dr \, d\theta \]
Evaluate the inner integral:
\[ \int_{r^2}^4 zr \, dz = r \left[ \frac{1}{2}z^2 \right]_{r^2}^4 = \frac{r}{2}(16 - r^4) = 8r - \frac{1}{2}r^5 \]
Evaluate the middle integral:
\[ \int_0^2 \left( 8r - \frac{1}{2}r^5 \right) \, dr = \left[ 4r^2 - \frac{1}{12}r^6 \right]_0^2 = 4(4) - \frac{64}{12} = 16 - \frac{16}{3} = \frac{48 - 16}{3} = \frac{32}{3} \]
Evaluate the outer integral:
\[ \int_0^{2\pi} \frac{32}{3} \, d\theta = \frac{32}{3}(2\pi) = \frac{64\pi}{3} \]
\textbf{Final Answers:} 
$A = 2\pi$, $B = 2$, $C = r^2$, Evaluated Value = $\frac{64\pi}{3}$

---

\subsection*{Problem 11}
\textbf{Question:} Evaluate $\iiint_E (x + y + z) \, dV$, where $E$ is the solid in the first octant that lies under the paraboloid $z = 16 - x^2 - y^2$.

\textbf{Solution:}
Convert to cylindrical coordinates:
Paraboloid boundary: $z = 16 - r^2$.
First octant constraints: $x \geq 0 \implies 0 \leq \theta \leq \frac{\pi}{2}$, and $z \geq 0$.
The base is in the $xy$-plane ($z=0$), so setting $0 = 16 - r^2 \implies r = 4$. Limits: $0 \leq r \leq 4$.
Integrand: $x + y + z = r\cos(\theta) + r\sin(\theta) + z$.
Volume differential: $r \, dz \, dr \, d\theta$.
Integral setup:
\[ \int_0^{\pi/2} \int_0^4 \int_0^{16-r^2} (r\cos(\theta) + r\sin(\theta) + z)r \, dz \, dr \, d\theta \]
Distribute $r$:
\[ \int_0^{16-r^2} \left( r^2(\cos(\theta)+\sin(\theta)) + zr \right) \, dz \]
Inner integral:
\[ \left[ z r^2(\cos(\theta)+\sin(\theta)) + \frac{z^2}{2}r \right]_0^{16-r^2} \]
\[ = (16-r^2)r^2(\cos(\theta)+\sin(\theta)) + \frac{1}{2}r(16-r^2)^2 \]
Middle integral:
Split into two parts: $\int_0^4 \dots \, dr$
Part 1: $\int_0^4 (16r^2 - r^4)(\cos(\theta)+\sin(\theta)) \, dr = (\cos(\theta)+\sin(\theta)) \left[ \frac{16r^3}{3} - \frac{r^5}{5} \right]_0^4$
\[ = (\cos(\theta)+\sin(\theta)) \left( \frac{1024}{3} - \frac{1024}{5} \right) = \frac{2048}{15}(\cos(\theta)+\sin(\theta)) \]
Part 2: $\int_0^4 \frac{1}{2}r(16-r^2)^2 \, dr$. Use u-sub: $u = 16-r^2, du = -2r \, dr \implies \frac{1}{2}r \, dr = -\frac{1}{4}du$.
Limits: $r=0 \to u=16$, $r=4 \to u=0$.
\[ \int_{16}^0 -\frac{1}{4}u^2 \, du = \frac{1}{4} \int_0^{16} u^2 \, du = \frac{1}{4} \left[ \frac{u^3}{3} \right]_0^{16} = \frac{4096}{12} = \frac{1024}{3} \]
Add Part 1 and 2:
\[ \text{Expression} = \frac{2048}{15}(\cos(\theta)+\sin(\theta)) + \frac{1024}{3} \]
Outer integral (integrate w.r.t $\theta$ from $0$ to $\pi/2$):
\[ \int_0^{\pi/2} \left( \frac{2048}{15}(\cos(\theta)+\sin(\theta)) + \frac{1024}{3} \right) \, d\theta \]
\[ = \frac{2048}{15} \left[ \sin(\theta) - \cos(\theta) \right]_0^{\pi/2} + \frac{1024}{3} [\theta]_0^{\pi/2} \]
\[ = \frac{2048}{15} ((1 - 0) - (0 - 1)) + \frac{1024}{3} \left(\frac{\pi}{2}\right) = \frac{2048}{15}(2) + \frac{512\pi}{3} = \frac{4096}{15} + \frac{512\pi}{3} \]
\textbf{Final Answer:} $\frac{4096}{15} + \frac{512\pi}{3}$

---

\subsection*{Problem 12}
\textbf{Question:} Evaluate $\iiint_E (x - y) \, dV$, where $E$ lies between the cylinders $x^2 + y^2 = 1$ and $x^2 + y^2 = 25$, above the $xy$-plane, and below the plane $z = y + 5$.

\textbf{Solution:}
Convert boundaries:
Cylinders: $r^2 = 1 \implies r=1$ and $r^2 = 25 \implies r=5$.
$xy$-plane: $z = 0$.
Upper plane: $z = r\sin(\theta) + 5$.
Angle limits: full circle, $0 \leq \theta \leq 2\pi$.
Integrand: $x - y = r\cos(\theta) - r\sin(\theta)$.
Integral setup:
\[ \int_0^{2\pi} \int_1^5 \int_0^{r\sin(\theta)+5} (r\cos(\theta) - r\sin(\theta))r \, dz \, dr \, d\theta \]
Inner integral:
\[ \int_0^{r\sin(\theta)+5} r^2(\cos(\theta) - \sin(\theta)) \, dz = r^2(\cos(\theta) - \sin(\theta))(r\sin(\theta) + 5) \]
Expand:
\[ = r^3(\cos(\theta)\sin(\theta) - \sin^2(\theta)) + 5r^2(\cos(\theta) - \sin(\theta)) \]
To make things easier, we can swap the $r$ and $\theta$ integrals (Fubini's Theorem) and integrate over $\theta$ first because many terms will cancel to zero over a full period $[0, 2\pi]$:
$\int_0^{2\pi} \cos(\theta)\sin(\theta) \, d\theta = 0$
$\int_0^{2\pi} \cos(\theta) \, d\theta = 0$
$\int_0^{2\pi} \sin(\theta) \, d\theta = 0$
$\int_0^{2\pi} \sin^2(\theta) \, d\theta = \pi$
So, integrating the expanded expression over $\theta$:
\[ \int_0^{2\pi} \left[ r^3(\cos(\theta)\sin(\theta) - \sin^2(\theta)) + 5r^2(\cos(\theta) - \sin(\theta)) \right] d\theta = r^3(0 - \pi) + 5r^2(0) = -\pi r^3 \]
Now perform the middle integral (which is now the last step) over $r$ from $1$ to $5$:
\[ \int_1^5 -\pi r^3 \, dr = -\pi \left[ \frac{r^4}{4} \right]_1^5 = -\pi \left( \frac{625 - 1}{4} \right) = -\pi \left( \frac{624}{4} \right) = -156\pi \]
\textbf{Final Answer:} $-156\pi$

---

\subsection*{Problem 13}
\textbf{Question:} Evaluate $\iiint_E x^2 \, dV$, where $E$ is the solid that lies within the cylinder $x^2 + y^2 = 4$, above the plane $z = 0$, and below the cone $z^2 = 16x^2 + 16y^2$.

\textbf{Solution:}
Convert boundaries:
Cylinder: $r^2 = 4 \implies r = 2$.
Cone: $z^2 = 16(x^2 + y^2) = 16r^2$. Since it's above $z=0$, $z = 4r$.
Limits: $0 \leq z \leq 4r$, $0 \leq r \leq 2$, $0 \leq \theta \leq 2\pi$.
Integrand: $x^2 = (r\cos(\theta))^2 = r^2\cos^2(\theta)$.
Integral setup:
\[ \int_0^{2\pi} \int_0^2 \int_0^{4r} (r^2\cos^2(\theta))r \, dz \, dr \, d\theta = \int_0^{2\pi} \int_0^2 \int_0^{4r} r^3\cos^2(\theta) \, dz \, dr \, d\theta \]
Inner integral:
\[ \int_0^{4r} r^3\cos^2(\theta) \, dz = r^3\cos^2(\theta) [z]_0^{4r} = 4r^4\cos^2(\theta) \]
Middle integral:
\[ \int_0^2 4r^4\cos^2(\theta) \, dr = 4\cos^2(\theta) \left[ \frac{r^5}{5} \right]_0^2 = 4\cos^2(\theta) \left(\frac{32}{5}\right) = \frac{128}{5}\cos^2(\theta) \]
Outer integral:
\[ \int_0^{2\pi} \frac{128}{5}\cos^2(\theta) \, d\theta = \frac{128}{5} \int_0^{2\pi} \frac{1 + \cos(2\theta)}{2} \, d\theta = \frac{64}{5} \left[ \theta + \frac{\sin(2\theta)}{2} \right]_0^{2\pi} \]
\[ = \frac{64}{5} (2\pi + 0) = \frac{128\pi}{5} \]
\textbf{Final Answer:} $\frac{128\pi}{5}$

---

\subsection*{Problem 14}
\textbf{Question:} Find the volume of the solid that lies within both the cylinder $x^2 + y^2 = 4$ and the sphere $x^2 + y^2 + z^2 = 49$.

\textbf{Solution:}
Volume implies we are evaluating $\iiint 1 \, dV$.
Cylinder bounds: $r^2 = 4 \implies 0 \leq r \leq 2$.
Sphere bounds for $z$: $r^2 + z^2 = 49 \implies z^2 = 49 - r^2 \implies -\sqrt{49 - r^2} \leq z \leq \sqrt{49 - r^2}$.
Angle limits: $0 \leq \theta \leq 2\pi$.
Integral setup:
\[ \int_0^{2\pi} \int_0^2 \int_{-\sqrt{49-r^2}}^{\sqrt{49-r^2}} r \, dz \, dr \, d\theta \]
Inner integral:
\[ \int_{-\sqrt{49-r^2}}^{\sqrt{49-r^2}} r \, dz = r \left( \sqrt{49-r^2} - (-\sqrt{49-r^2}) \right) = 2r\sqrt{49-r^2} \]
Middle integral:
\[ \int_0^2 2r\sqrt{49-r^2} \, dr \]
Use u-substitution: $u = 49 - r^2, du = -2r \, dr \implies -du = 2r \, dr$.
Limits: $r=0 \to u=49$, $r=2 \to u=45$.
\[ \int_{49}^{45} -\sqrt{u} \, du = \int_{45}^{49} u^{1/2} \, du = \left[ \frac{2}{3}u^{3/2} \right]_{45}^{49} = \frac{2}{3}(49^{3/2} - 45^{3/2}) \]
Note that $49^{3/2} = (\sqrt{49})^3 = 7^3 = 343$. And $45^{3/2} = (\sqrt{45})^3 = (3\sqrt{5})^3 = 135\sqrt{5}$.
\[ = \frac{2}{3}(343 - 135\sqrt{5}) \]
Outer integral:
\[ \int_0^{2\pi} \frac{2}{3}(343 - 135\sqrt{5}) \, d\theta = \frac{4\pi}{3}(343 - 135\sqrt{5}) \]
\textbf{Final Answer:} $\frac{4\pi}{3}(343 - 135\sqrt{5})$

---

\subsection*{Problem 15}
\textbf{Question:} Find the volume of the solid enclosed by the cone $z = \sqrt{x^2 + y^2}$ and the sphere $x^2 + y^2 + z^2 = 8$.

\textbf{Solution:}
Convert boundaries:
Cone: $z = \sqrt{r^2} = r$.
Sphere: $r^2 + z^2 = 8 \implies z = \sqrt{8 - r^2}$ (upper hemisphere limits since it's bounded by the upward opening cone).
Find intersection to get $r$-limits: $r = \sqrt{8 - r^2} \implies r^2 = 8 - r^2 \implies 2r^2 = 8 \implies r^2 = 4 \implies r = 2$.
Integral setup:
\[ V = \int_0^{2\pi} \int_0^2 \int_r^{\sqrt{8-r^2}} r \, dz \, dr \, d\theta \]
Inner integral:
\[ \int_r^{\sqrt{8-r^2}} r \, dz = r(\sqrt{8-r^2} - r) = r\sqrt{8-r^2} - r^2 \]
Middle integral:
\[ \int_0^2 (r\sqrt{8-r^2} - r^2) \, dr \]
Split into two:
Part 1: $\int_0^2 r\sqrt{8-r^2} \, dr$. Let $u = 8-r^2$, $du = -2r \, dr$. Limits: $8$ to $4$.
\[ \int_8^4 -\frac{1}{2}u^{1/2} \, du = \frac{1}{2} \int_4^8 u^{1/2} \, du = \frac{1}{2} \left[ \frac{2}{3}u^{3/2} \right]_4^8 = \frac{1}{3}(8^{3/2} - 4^{3/2}) = \frac{1}{3}(16\sqrt{2} - 8) \]
Part 2: $\int_0^2 r^2 \, dr = \left[ \frac{r^3}{3} \right]_0^2 = \frac{8}{3}$.
Combine:
\[ \frac{16\sqrt{2} - 8}{3} - \frac{8}{3} = \frac{16\sqrt{2} - 16}{3} \]
Outer integral:
\[ \int_0^{2\pi} \frac{16\sqrt{2} - 16}{3} \, d\theta = 2\pi \left( \frac{16(\sqrt{2} - 1)}{3} \right) = \frac{32\pi}{3}(\sqrt{2} - 1) \]
\textbf{Final Answer:} $\frac{32\pi}{3}(\sqrt{2} - 1)$

---

\subsection*{Problem 16}
\textbf{Question:} Consider the volume of the solid enclosed by the cone $z = \sqrt{x^2 + y^2}$ and the sphere $x^2 + y^2 + z^2 = 72$. Write an integral: $\int_0^A \int_0^B \int_r^C r \, dz \, dr \, d\theta$. Find $A, B, C$ and the volume.

\textbf{Solution:}
Convert boundaries:
Cone: $z = r$.
Sphere: $z = \sqrt{72 - r^2}$. (This is the upper bound $C$).
Intersection: $r = \sqrt{72 - r^2} \implies 2r^2 = 72 \implies r^2 = 36 \implies r = 6$. So $0 \leq r \leq 6$, making $B = 6$.
Full rotation implies $A = 2\pi$.
Integral setup:
\[ V = \int_0^{2\pi} \int_0^6 \int_r^{\sqrt{72-r^2}} r \, dz \, dr \, d\theta \]
Inner integral:
\[ \int_r^{\sqrt{72-r^2}} r \, dz = r\sqrt{72-r^2} - r^2 \]
Middle integral:
\[ \int_0^6 (r\sqrt{72-r^2} - r^2) \, dr \]
Part 1: $\int_0^6 r\sqrt{72-r^2} \, dr$. Let $u=72-r^2$, $du=-2r \, dr$. Limits $72$ to $36$.
\[ = \frac{1}{2} \int_{36}^{72} u^{1/2} \, du = \frac{1}{3} [u^{3/2}]_{36}^{72} = \frac{1}{3}(72\sqrt{72} - 36\sqrt{36}) = \frac{1}{3}(432\sqrt{2} - 216) = 144\sqrt{2} - 72 \]
Part 2: $\int_0^6 r^2 \, dr = \left[\frac{r^3}{3}\right]_0^6 = \frac{216}{3} = 72$.
Subtract Part 2 from Part 1:
\[ 144\sqrt{2} - 72 - 72 = 144\sqrt{2} - 144 = 144(\sqrt{2} - 1) \]
Outer integral:
\[ \int_0^{2\pi} 144(\sqrt{2} - 1) \, d\theta = 288\pi(\sqrt{2} - 1) \]
\textbf{Final Answers:}
$A = 2\pi$, $B = 6$, $C = \sqrt{72-r^2}$, Volume = $288\pi(\sqrt{2} - 1)$

\newpage
\section{Part 3: In-Depth Analysis of Problems and Techniques}

\subsection{A) Problem Types and General Approach}
The homework problems can be categorized into four distinct archetypes:

\textbf{1. Coordinate Conversion Problems (Q1-Q3)}
\textit{Strategy:} Use direct substitution with the transformation formulas ($x = r\cos(\theta)$, etc.). The main trap is properly identifying the correct quadrant for $\theta$ when inverse tangent yields multiple possibilities.

\textbf{2. Surface Identification and Translation (Q4-Q6)}
\textit{Strategy:} These require algebraic manipulation to group variables naturally suited for cylindrical form. You gather $x^2$ and $y^2$ terms to form $r^2$. Recognizing standard forms (like spheres, cones, or paraboloids) provides a physical intuition check.

\textbf{3. Boundary Intersection and Setup Problems (Q8, Q10, Q16)}
\textit{Strategy:} The hardest part of calculus in 3D is setting up the bounds. The general strategy is:
1. Write the "floor" and "ceiling" functions in terms of $r$ to find the $z$-bounds.
2. Set the floor and ceiling equal to each other ($z_{\text{top}} = z_{\text{bottom}}$) to find the shadow region in the $xy$-plane, which dictates the $r$-bounds.
3. Determine the rotational sweep to find $\theta$-bounds.

\textbf{4. Full Iterated Evaluation Problems (Q7, Q9, Q11-Q15)}
\textit{Strategy:} These require the mechanical application of integration. A common, powerful trick applied here is separating the integrals or altering the order of integration (Fubini's Theorem) when variables are independent or expressions evaluate trivially over a full period (as seen in Q12).

\subsection{B) Key Algebraic and Calculus Manipulations}

\textbf{The Fubini Simplification Trick (Seen in Q12):}
In Problem 12, the expanded integrand was complicated: $r^3\cos(\theta)\sin(\theta) - r^3\sin^2(\theta) \dots$ Integrating this with respect to $r$ first would have created messy fractions. By recognizing that we integrate over a full period $[0, 2\pi]$, we swapped the order to integrate with respect to $\theta$ first. Since $\int_0^{2\pi} \sin(\theta) \, d\theta = 0$, three of the four massive terms collapsed to zero instantly.

\textbf{U-Substitution for Fractional Powers (Seen in Q14, Q15, Q16):}
When finding the volume of spheres, you inevitably get integrals of the form $\int r\sqrt{R^2 - r^2} \, dr$. The algebraic trick here is recognizing that the $r$ provided by the Jacobian determinant ($dV = r \, dz \, dr \, d\theta$) is exactly what is needed to execute a $u$-substitution. Without that extra $r$, the integral would require complex trigonometric substitution.

\textbf{Double Angle Identity for Powers of Cosine (Seen in Q5b, Q13):}
In problem 13, evaluating $\int \cos^2(\theta) \, d\theta$ was required. You cannot use power rule here; you must use the identity $\cos^2(\theta) = \frac{1 + \cos(2\theta)}{2}$.

\newpage
\section{Part 4: "Cheatsheet" and Tips for Success}

\subsection{Crucial Formulas Summary}
\begin{itemize}
    \item \textbf{Cylindrical to Rectangular:} $x = r \cos(\theta)$, $y = r \sin(\theta)$, $z = z$
    \item \textbf{Rectangular to Cylindrical:} $r^2 = x^2 + y^2$, $\tan(\theta) = \frac{y}{x}$ (check quadrant!)
    \item \textbf{The Essential Jacobian Element:} $dV = r \, dz \, dr \, d\theta$
\end{itemize}

\subsection{Quick Recognition: Standard Surfaces in Cylindrical}
\begin{itemize}
    \item \textbf{Cylinder:} $r = c$ (e.g., $r = 5$)
    \item \textbf{Cone:} $z = c \cdot r$ (e.g., $z = r$)
    \item \textbf{Paraboloid:} $z = c \cdot r^2$ (e.g., $z = r^2$)
    \item \textbf{Sphere:} $r^2 + z^2 = R^2$
    \item \textbf{Horizontal Plane:} $z = c$
    \item \textbf{Vertical Plane:} $\theta = c$ (e.g., $\theta = \frac{\pi}{4}$)
\end{itemize}

\subsection{Shortcuts and "Cheats"}
\begin{itemize}
    \item \textbf{The Odd Function Symmetry Cheat:} If your bounds for $\theta$ are a full circle $[0, 2\pi]$ and your integrand consists of a single $\sin(\theta)$ or $\cos(\theta)$ to an odd power (like in Problem 12), the integral evaluates to zero. Don't waste time integrating the rest of the variables!
    \item \textbf{Independence of Variables:} If your bounds are all constants AND the function can be factored into $f(r)g(\theta)h(z)$, you can separate the integral into a multiplication of three single integrals: $(\int h(z)dz) \cdot (\int f(r)dr) \cdot (\int g(\theta)d\theta)$.
\end{itemize}

\subsection{Common Pitfalls to Watch Out For}
\begin{itemize}
    \item \textbf{Forgetting the Jacobian $r$:} The \#1 mistake students make is writing $dV = dz \, dr \, d\theta$. It \textit{must} be $r \, dz \, dr \, d\theta$.
    \item \textbf{Quadrant Errors with $\tan^{-1}$:} In Problem 3b, $y/x$ is negative. Your calculator might spit out a negative angle in quadrant IV, but you must ensure it matches the actual $(x,y)$ coordinate plotted physically.
    \item \textbf{Squaring Errors:} When substituting $x^2 + y^2$ with $r^2$, remember that $(x+y)^2 \neq r^2$. It only works for the sum of squares.
\end{itemize}

\newpage
\section{Part 5: Conceptual Synthesis and The "Big Picture"}

\subsection{A) Thematic Connections}
The core theme of this chapter is \textbf{Geometric Transformation to Exploit Symmetry}. Mathematics relies heavily on choosing the right perspective. Just as a matrix diagonalization changes a basis to make a complicated linear transformation simple, cylindrical coordinates change the perspective of the grid. Instead of forcing curved boundaries to fit into rectangular boxes (which causes aggressive square roots and inverse trig functions), we bend our measurement axes to perfectly wrap around the object. This theme is universal in higher mathematics: adapt your tools to the shape of the problem, not vice versa.

\subsection{B) Forward and Backward Links}
\textbf{Backward Links:} This topic is a direct synthesis of two previous concepts. It takes 2D polar integration from Calculus II and marries it to 3D rectangular triple integrals from early Calculus III. It proves that the logic used to derive $dA = r \, dr \, d\theta$ scales perfectly into higher dimensions.

\textbf{Forward Links:} Mastery of cylindrical coordinates is absolutely required for the upcoming jump into Spherical Coordinates (which involves two angles and one radius). Beyond Calculus III, this forms the foundation for Vector Calculus (specifically Stokes' Theorem and the Divergence Theorem). When calculating flux across a cylindrical pipe in fluid dynamics, these coordinate conversions are not just helpful; they are mandatory. 

\newpage
\section{Part 6: Real-World Application and Modeling}

\subsection{A) Concrete Scenarios (Finance/Stats Focus)}
While direct physical volume calculations map easily to engineering, triple integrals over cylindrical/spherical geometries also hold profound applications in quantitative finance, stochastic calculus, and statistical modeling. 

\textbf{1. Pricing Basket Options with Correlated Assets:} In exotic derivative pricing, a "basket option" payout depends on the performance of multiple underlying assets. If asset X and asset Y possess a correlation structure that yields a circularly symmetric joint probability distribution (like a bivariate normal distribution with equal variances), determining the expected payoff requires integrating over a cylindrical region in 3D space, where the $z$-axis represents the payout function.

\textbf{2. Trivariate Probability Density Functions:} In risk management (Value at Risk - VaR modeling), quants analyze the joint behavior of three factors (e.g., Interest Rates, Inflation, and Equity Index returns). If these factors follow a spherically or cylindrically symmetric joint density distribution, computing the probability of a market crash (the volume under the probability surface) demands cylindrical or spherical coordinate integration.

\subsection{B) Model Problem Setup: The Probability of Joint Default}
\textbf{Scenario Description:} A financial institution models the default risk of two correlated tech companies (Company X and Company Y). Their standardized return parameters, $x$ and $y$, have a joint circular normal distribution. The severity of the loss, $z$, depends quadratically on their joint variance. 

\textbf{Model Formulation:}
Let the joint probability density function of the returns be a standard bivariate normal distribution, which has circular symmetry:
\[ f(x, y) = \frac{1}{2\pi} e^{-\frac{x^2 + y^2}{2}} \]
Let the loss function (the $z$-dimension) be defined by the paraboloid:
\[ L(x, y) = x^2 + y^2 \]
A "Severe Default Event" is defined as a scenario where the joint squared returns are less than or equal to $R^2$, and the loss severity is above a critical threshold $L_{min}$ but below $L_{max}$.

To find the Expected Value of the loss strictly within this default scenario, we set up a triple integral. In rectangular form, it's messy. In cylindrical coordinates ($r^2 = x^2 + y^2$), the density function simplifies to $f(r) = \frac{1}{2\pi}e^{-r^2/2}$ and the loss limits become $L_{min} \leq z \leq L_{max}$.

\textbf{The Integral Model:}
\[ \text{Expected Loss} = \int_0^{2\pi} \int_0^R \int_{L_{min}}^{L_{max}} \left( z \cdot \frac{1}{2\pi}e^{-r^2/2} \right) r \, dz \, dr \, d\theta \]
This model instantly converts an impossible statistical risk problem into a trivial separation-of-variables integral!

\newpage
\section{Part 7: Common Variations and Untested Concepts}

The homework assignment focused heavily on pure geometric volumes. However, standard Calculus III curricula strongly emphasize \textbf{Density, Mass, and Center of Mass}. 

\textbf{Missing Concept: Mass of a Solid with Variable Density}
If a solid has a non-uniform density $\rho(x,y,z)$, its mass is not just its volume. You must integrate the density function over the volume:
\[ \text{Mass } m = \iiint_E \rho(r, \theta, z) \, dV \]

\textbf{Worked Example:}
Find the mass of a solid cylinder $E$ defined by $0 \leq r \leq 2$, $0 \leq z \leq 5$, whose density is proportional to the distance from the central axis. 
\textit{Solution:}
The distance from the $z$-axis is $\sqrt{x^2+y^2} = r$. Thus, density $\rho(x,y,z) = K \cdot r$ where $K$ is a constant.
\[ m = \int_0^{2\pi} \int_0^2 \int_0^5 (K \cdot r) r \, dz \, dr \, d\theta \]
\[ = K \int_0^{2\pi} d\theta \int_0^2 r^2 \, dr \int_0^5 dz \]
\[ = K (2\pi) \left[ \frac{8}{3} \right] (5) = \frac{80\pi K}{3} \]

\newpage
\section{Part 8: Advanced Diagnostic Testing: "Find the Flaw"}

\textbf{Instructions for the Student:} Below are 5 newly generated problems. The provided solutions contain exactly one critical flaw. Identify the specific error, explain why it is wrong in one sentence, and provide the correct calculation.

---
\textbf{Flawed Problem 1:} 
Evaluate the volume of the region bounded by $z = x^2 + y^2$ and $z = 9$.
\textit{Flawed Solution:} 
Intersection occurs at $r^2 = 9 \implies r = 3$. 
$V = \int_0^{2\pi} \int_0^3 \int_{r^2}^9 1 \, dz \, dr \, d\theta$.
Inner: $[z]_{r^2}^9 = 9 - r^2$. 
Middle: $\int_0^3 (9 - r^2) \, dr = [9r - \frac{r^3}{3}]_0^3 = 27 - 9 = 18$.
Outer: $\int_0^{2\pi} 18 \, d\theta = 36\pi$.

---
\textbf{Flawed Problem 2:}
Convert the equation $x^2 - y^2 + z^2 = 1$ to cylindrical coordinates.
\textit{Flawed Solution:}
We know that in cylindrical coordinates, the sum of squares translates directly to the radius. So $x^2 - y^2$ becomes $r^2$.
Therefore, the equation is $r^2 + z^2 = 1$.

---
\textbf{Flawed Problem 3:}
Evaluate $\iiint_E z \, dV$ where $E$ is the region under $z=4$ and above the $xy$-plane, in the first octant only, inside the cylinder $x^2+y^2=1$.
\textit{Flawed Solution:}
First octant means $x>0, y>0, z>0$. 
$\int_0^{\pi} \int_0^1 \int_0^4 z \cdot r \, dz \, dr \, d\theta$.
Inner: $[z^2/2]_0^4 = 8$.
Middle: $\int_0^1 8r \, dr = [4r^2]_0^1 = 4$.
Outer: $\int_0^{\pi} 4 \, d\theta = 4\pi$.

---
\textbf{Flawed Problem 4:}
Find the volume between the cone $z = 2r$ and the plane $z=0$, inside the cylinder $r=4$.
\textit{Flawed Solution:}
The region lies above the $xy$-plane and below the cone.
$V = \int_0^{2\pi} \int_0^4 \int_{2r}^0 r \, dz \, dr \, d\theta$.
Inner: $\int_{2r}^0 r \, dz = r[0 - 2r] = -2r^2$.
Middle: $\int_0^4 -2r^2 \, dr = -128/3$.
Outer: evaluates to $-256\pi/3$. Volume is $256\pi/3$ by taking absolute value.

---
\textbf{Flawed Problem 5:}
Set up the integral to find the volume inside $x^2+y^2=9$ and between $z=-1$ and $z=2$.
\textit{Flawed Solution:}
Cylinder is $r=3$.
$V = \int_{-1}^{2} \int_{0}^{3} \int_{0}^{2\pi} 1 \, d\theta \, dr \, dz$.
Inner: $2\pi$. 
Middle: $\int_0^3 2\pi \, dr = 6\pi$.
Outer: $\int_{-1}^2 6\pi \, dz = 6\pi(3) = 18\pi$.

---
\subsection*{Solutions to "Find the Flaw"}
\textit{(To be checked by the student, provided here for the tutor's reference)}

\textbf{Flaw 1 Correction:} The flaw is omitting the Jacobian $r$ in the integrand. The volume element must be $r \, dz \, dr \, d\theta$. Correct calculation yields $V = \frac{81\pi}{2}$.
\textbf{Flaw 2 Correction:} The flaw is incorrectly assuming $x^2 - y^2 = r^2$. Only $x^2 + y^2 = r^2$; $x^2 - y^2$ becomes $r^2\cos(2\theta)$. Correct equation is $r^2\cos(2\theta) + z^2 = 1$.
\textbf{Flaw 3 Correction:} The flaw is using bounds $0$ to $\pi$ for the first octant. The first octant restricts $\theta$ to $[0, \pi/2]$. Correct calculation yields $\pi$.
\textbf{Flaw 4 Correction:} The flaw is reversing the $z$ bounds, placing $2r$ on bottom and $0$ on top, which forces a negative volume. The solid between the cone and plane inside $r=4$ is technically bounded above by $z=2r$ and below by $z=0$, so bounds must be $\int_0^{2r}$. 
\textbf{Flaw 5 Correction:} The flaw is again omitting the Jacobian $r$. While the integral bounds and order $d\theta \, dr \, dz$ are valid by Fubini's, the integrand must be $r$, not $1$. Correct calculation yields $27\pi$.

\end{document}